% Secondary figure: `small-multiples` (design-axes.md axis 10; P1's bars,
% P4's 2x2 panels). One small panel per condition, arms as bars, shared
% y-axis. Tier B from data/opening.dat; data/opening-multiples.tex is the
% groupplot body written by make_paper_data.py --opening.
%
% source-data: runs/aggregate__*.json -> paper/data/opening.dat
% generator:   scripts/make_paper_data.py --opening
\begin{acwidefigure}
  \centering
  \begin{tikzpicture}
    \begin{groupplot}[group style={group size=\OPncondsnum\space by 1, horizontal sep=0.35cm,
                                   ylabels at=edge left, yticklabels at=edge left},
      width=\OPpanelwidth, height=4.0cm,
      ymin=\OPymin, ymax=\OPymax, xtick=\empty, enlarge x limits=0.35,
      ymajorgrids, grid style={ACMuted!25}, axis line style={ACInk!70},
      tick label style={font=\tiny}, label style={font=\scriptsize},
      title style={font=\tiny,yshift=-2pt,align=center,text width=\OPpanelwidth}, ylabel={\OPmetric\ \OPdirection},
      legend style={font=\tiny,draw=none,fill=none,at={(0,1.18)},anchor=south west,legend columns=-1},
    ]
      \input{data/opening-multiples}
    \end{groupplot}
  \end{tikzpicture}
  \caption{\textbf{\slot{name what the panels compare}.} One panel per
    condition; bars are means over repetitions, in the arm order of the
    legend. \slot{one sentence stating what to notice across panels}.}
  \label{fig:distribution}
\end{acwidefigure}
